\documentclass[11pt,english]{article}

\usepackage[T1]{fontenc}
\usepackage[sc]{mathpazo}

\usepackage[a4paper,
    top=1in,
    bottom=1in,
    left=1in,
    right=1in]{geometry}

\usepackage{amsmath,amsthm,amssymb,mathtools}
\usepackage{graphicx}
\usepackage{booktabs,tabularx,longtable}
\usepackage{enumitem}
\usepackage[dvipsnames]{xcolor}
\usepackage{setspace}
\usepackage{fancyhdr}
\usepackage{hyperref}

\onehalfspacing
\setlength{\parskip}{1em}

\hypersetup{
    colorlinks=true,
    linkcolor=black,
    citecolor=black,
    urlcolor=blue
}

\begin{document}

We modeling penalty corners in field hockey as a game between the offense and defense. At each corner, the offense chooses a setup and strategy with the objective of scoring, while the defense chooses a strategy to prevent a goal. Corners are generated by fouls committed inside the circle, so the number of corners in a match is not known in advance. In particular, any corner may be the final corner of the match, but teams do not know ex ante when the final corner will occur.

We therefore model corners as an indefinite horizon zero sum repeated game in which each corner constitutes a stage game. Each stage game has two parts: teams first observe the current state of the match, and then simultaneously choose offensive and defensive strategies.
\[
\text{For each corner } t:
\qquad
\left\{
\begin{aligned}
S_t
&=
(H_t,\ \text{score}_t,\ \text{time}_t,\ F_t), \\[4pt]
\sigma^O(\cdot \mid S_t)
&\in \Delta(\mathcal{A}^O),
\qquad
\sigma^D(\cdot \mid S_t)
\in \Delta(\mathcal{A}^D), \\[4pt]
A_t^O
&\sim \sigma^O(\cdot \mid S_t),
\qquad
A_t^D
\sim \sigma^D(\cdot \mid S_t), \\[4pt]
Y_t
&\in
\{\text{Goal},\ \text{No Goal}\}, \\[4pt]
H_{t+1}
&=
(H_t,\ A_t^O,\ A_t^D,\ Y_t), \\[4pt]
S_{t+1}
&=
(H_{t+1},\ \text{score}_{t+1},\ \text{time}_{t+1},\ F_{t+1}).
\end{aligned}
\right.
\]
where
\begin{itemize}
    \item $S_t$ is the state of the game at corner $t$, composed of the relevant corner history, the game score, the game time, and the observed APC formation.

    \item $\mathcal{A}^O$ and $\mathcal{A}^D$ are the sets of possible offensive and defensive actions, respectively.

    \item $\sigma^O(\cdot \mid S_t)$ and $\sigma^D(\cdot \mid S_t)$ are the offensive and defensive strategies conditional on the state $S_t$. Each strategy specifies a probability distribution over the relevant action set.

    \item $A_t^O$ and $A_t^D$ are the realized offensive and defensive actions drawn from these strategies.

    \item $Y_t$ is the realized outcome of the corner.

    \item $S_{t+1}$ is the state at the next corner. Thus, after each corner, the state is updated to incorporate the new information generated by the current corner.
\end{itemize}

\begin{center}
\fbox{
\begin{minipage}{0.92\textwidth}


\textbf{Stage Game}

\vspace{0.4em}

\textbf{Part 1: Common information is revealed.}  
Once teams know that a penalty corner will occur, they observe the relevant game state, including the time remaining, score, APC initial setup, and other publicly available information.

\vspace{0.6em}

\textbf{Part 2: Strategies are chosen simultaneously.}  
Conditional on this information, the offense and defense simultaneously choose their strategies. The resulting outcome is then realized (e.g., goal or no goal).

\end{minipage}
}
\end{center} 
We know that a mixed strategy Nash equilibrium exists in this game for each state $S_t$ by Nash's existence theorem.  Thus we can calculate best response correspondences for each team. 


\noindent \textbf{Proof of Concept/Simple Example:}

Consider some state $S_t$.
\[
\text{The offense has strategies: }
\sigma_O(\cdot \mid S_t)
\text{ where their action } A_t^O
\text{ will be drawn from.}
\]
\[
\text{The defense has strategies: }
\sigma_D(\cdot \mid S_t)
\text{ where their action } A_t^D
\text{ is drawn from.}
\]

The offense maximizes:

\[
u^O(A_t^O, A_t^D, S_t)
=
P(Y_t = \text{Goal} \mid A_t^O, A_t^D, S_t)
\]

Since the game is zero sum, the defense gets:

\[
u^D(A_t^O, A_t^D, S_t)
=
- u^O
\]

Using existing data, we have some payoff matrix that we can populate for each $S_t$:

\[
\begin{array}{c|ccccc}
    & A_1^O & A_2^O & A_3^O & \cdots & A_n^O \\ \hline
A_1^D & & & & & \\
A_2^D & & & & & \\
A_3^D & & & & & \\
\vdots & & & & & \\
A_h^D & & & & &
\end{array}
\]

This is what we need to calculate with data. Each entry is just

\[
P(\text{Goal} \mid A_t^D, A_t^O, S_t).
\]

For example, with two strategies, we might have the following
payoff matrix:

$$
\begin{array}{c|cc}
    & A_1^O & A_2^O \\ \hline
A_1^D & 0.5 & 0.5 \\
A_2^D & 0.75 & 0.25
\end{array}
$$

which would represent the following game, where offense maximizes
the probability of a goal and defense minimizes it. Equivalently,
if defense utility is defined as the negative probability of a goal,
both players maximize their own utility:

$$
\begin{array}{c|cc}
    & A_1^O & A_2^O \\ \hline
A_1^D & \fcolorbox{green}{white}{$(-.5,.5)$} & (-.5,.5) \\
A_2^D & (-.75,.75) & (-.25,.25)
\end{array}
$$

The upper-left cell is a Nash equilibrium. Given that the offense
plays $A_1^O$, the defense prefers $A_1^D$, since

$$
-.5 > -.75.
$$

Given that the defense plays $A_1^D$, the offense is indifferent
between $A_1^O$ and $A_2^O$, since both give payoff $.5$.
Thus neither player has a profitable unilateral deviation from
$(A_1^D,A_1^O)$.


We could also have a game like this:

\[
\begin{array}{c|cc}
    & A_1^O & A_2^O \\ \hline
A_1^D & (-.25,.25) & (-.75,.75) \\
A_2^D & (-.75,.75) & (-.25,.25)
\end{array}
\]

Here, there is no pure strategy Nash equilibrium.

Let $p = \Pr(A_1^O)$ and $q = \Pr(A_1^D)$. We look for indifference
point:
\[
0.75p + 0.25(1-p) = .25p + 0.75(1-p)
\]
\[
p = 1/2
\]

and
\[
0.25q + 0.75(1-q) = 0.75q + 0.25(1-q)
\]
\[
q = 1/2
\]

Thus the Nash is $p = 1/2$, $q = 1/2$.

When we have a larger matrix, we can eliminate any dominated strategies and then solve for a Nash equilibrium. The resulting equilibrium will specify the probability with which each team should play each strategy in its equilibrium support.



Open Questions:
\begin{itemize}
    \item How should history be populated? 
    \begin{itemize}
        \item Could be related to just this game, to past games/scouting, to team practices with certain corners and so on. 
    \end{itemize}
    \item Are there other variables we can include to make $S_t$ more realistic? Anything we can remove? How granular can we make it? 
\end{itemize}



\end{document}